\chapter{上极限与下极限}
\label{chap:limsup-liminf}

对每个尾部取上确界和下确界，会得到一条递减列和一条递增列。它们的极限分别记住原数列长期可能达到的最高、最低水平，即使原数列本身发散也仍有意义。适当的子列能够达到这两个值，从而把确界估计与子列抽取联系起来。

\section{扩充实数与尾部上下确界}

\subsection{扩充实数系 \texorpdfstring{\([-\infty,+\infty]\)}{[-infinity,+infinity]}}


在集合 \(\R\) 中加入两个符号 \(-\infty,+\infty\)，得到\term{扩充实数系}
\[
 \overline{\R}=[-\infty,+\infty].
\]
规定
\[
 -\infty<x<+\infty\qquad(x\in\R).
\]
扩充实数是全序集，但不是域。表达式
\[
 (+\infty)+(-\infty),\qquad 0\cdot(+\infty),\qquad
 \frac{\infty}{\infty}
\]
不作定义，因为不同逼近方式会产生不同结果。只有不会造成歧义的运算才使用，例如 \(x+(+\infty)=+\infty\) 对有限 \(x\) 成立。

\begin{definition}
在 \(\overline{\R}\) 中，若数列趋于 \(+\infty\) 或 \(-\infty\)，也把相应符号称为其扩充实数极限。
\end{definition}

任意集合 \(E\subseteq\R\) 在扩充实数中都有上确界与下确界：若 \(E\) 在实数中无上界，规定 \(\sup E=+\infty\)；若 \(E=\varnothing\)，为统一尾部公式规定
\[
 \sup\varnothing=-\infty,\qquad
 \inf\varnothing=+\infty.
\]
本章尾部值集均非空，因此空集约定不会直接进入数列计算。

\subsection{尾部上确界与下确界}


对实数列 \((a_n)\)，定义扩充实数
\[
 \beta_n=\sup_{k\ge n}a_k,\qquad
 \alpha_n=\inf_{k\ge n}a_k.
\]
称它们为第 \(n\) 个尾部的上确界和下确界。

\begin{proposition}
\((\beta_n)\) 在 \(\overline{\R}\) 中单调递减，\((\alpha_n)\) 单调递增，并且
\[
 \alpha_n\le a_n\le\beta_n.
\]
若原列有界，则两列均为有限实数列并分别收敛。
\end{proposition}

\begin{proof}
尾部集合满足 \(A_{n+1}\subseteq A_n\)，所以其上确界不增、下确界不减。因 \(a_n\in A_n\)，中间不等式成立。有界时 \(\beta_n,\alpha_n\) 有共同有限界，单调有界收敛定理给出极限。
\end{proof}

\begin{example}
若 \(a_n=(-1)^n+1/n\)，则每个尾部仍含趋近 \(1\) 的偶数项和趋近 \(-1\) 的奇数项。直接计算尾部确界的细节依赖首个偶指标，但有
\[
 \beta_n\downarrow1,\qquad \alpha_n\uparrow-1.
\]
\end{example}

\section{\texorpdfstring{\(\limsup\)、\(\liminf\)}{limsup, liminf} 与单调尾列}


\begin{definition}
实数列 \((a_n)\) 的\term{上极限}与\term{下极限}定义为
\[
 \limsup_{n\to\infty}a_n
 :=\inf_{n\ge1}\sup_{k\ge n}a_k
 =\inf_n\beta_n,
\]
\[
 \liminf_{n\to\infty}a_n
 :=\sup_{n\ge1}\inf_{k\ge n}a_k
 =\sup_n\alpha_n.
\]
二者取值于 \(\overline{\R}\)。
\end{definition}

有界时，它们分别是单调列 \(\beta_n,\alpha_n\) 的实数极限。定义中的顺序不能交换：先对每个尾部取确界，再对尾部起点取下确界，表达“删去任意有限多项后仍可能达到的最高层次”。

上极限不是全体项的上确界。数列 \(a_1=100\)、\(a_n=0\)（\(n\ge2\)）的上确界为 \(100\)，上极限却是 \(0\)，因为有限个早期异常值最终会从尾部中消失。相反，只要某一高度在任意晚的尾部仍被项反复逼近，它就会影响上极限。

\begin{theorem}[基本序关系]
对任意实数列，
\[
 \liminf a_n\le\limsup a_n.
\]
若 \(a_n\le b_n\) 最终成立，则
\[
 \liminf a_n\le\liminf b_n,\qquad
 \limsup a_n\le\limsup b_n.
\]
\end{theorem}

\begin{proof}
对任意 \(m,n\)，取 \(k\ge\max\{m,n\}\)，有
\[
 \alpha_m\le a_k\le\beta_n.
\]
所以 \(\alpha_m\le\beta_n\)。先对 \(m\) 取上确界、再对 \(n\) 取下确界得到第一式。最终比较允许删去有限首项；在共同尾部上，尾部上下确界保持相应序关系，取极限即得第二组不等式。
\end{proof}

\begin{proposition}
\[
 \limsup(-a_n)=-\liminf a_n,\qquad
 \liminf(-a_n)=-\limsup a_n.
\]
对常数 \(c\in\R\)，
\[
 \limsup(a_n+c)=\limsup a_n+c,
\quad
 \liminf(a_n+c)=\liminf a_n+c.
\]
\end{proposition}

\begin{proof}
对每个 \(n\)，尾部集合在取负号后满足
\[
 \sup_{k\ge n}(-a_k)=-\inf_{k\ge n}a_k,
 \qquad
 \inf_{k\ge n}(-a_k)=-\sup_{k\ge n}a_k.
\]
令 \(n\to\infty\) 即得第一组等式。平移常数不改变尾部中各项的相对次序，并且
\[
 \sup_{k\ge n}(a_k+c)=\sup_{k\ge n}a_k+c,
 \qquad
 \inf_{k\ge n}(a_k+c)=\inf_{k\ge n}a_k+c.
\]
再令 \(n\to\infty\)，得到第二组等式。
\end{proof}

\begin{theorem}[上极限与下极限的阈值刻画]
设 \(s\in\R\)。则 \(\limsup a_n=s\) 当且仅当以下两项同时成立：
\begin{enumerate}
 \item 对每个 \(\varepsilon>0\)，最终有 \(a_n<s+\varepsilon\)；
 \item 对每个 \(\varepsilon>0\) 和每个 \(N\in\Npos\)，存在
 \(n\ge N\) 使 \(a_n>s-\varepsilon\)。
\end{enumerate}
对偶地，\(\liminf a_n=\ell\in\R\) 当且仅当：
\begin{enumerate}
 \item 对每个 \(\varepsilon>0\)，最终有 \(a_n>\ell-\varepsilon\)；
 \item 对每个 \(\varepsilon>0\) 和每个 \(N\in\Npos\)，存在
 \(n\ge N\) 使 \(a_n<\ell+\varepsilon\)。
\end{enumerate}
\end{theorem}

\begin{proof}
设尾部上确界为 \(\beta_N\)。若 \(\beta_N\downarrow s\)，对给定
\(\varepsilon>0\)，可取 \(N\) 使 \(\beta_N<s+\varepsilon\)，于是
\(n\ge N\) 时 \(a_n\le\beta_N<s+\varepsilon\)。另一方面，对任意
\(N\)，数 \(s-\varepsilon\) 不能是尾部的上界；否则
\(\beta_N\le s-\varepsilon<s\)，与 \(s=\inf_N\beta_N\) 矛盾。因此该尾部中存在 \(a_n>s-\varepsilon\)。

反之，第一项说明某个尾部上确界不超过 \(s+\varepsilon\)，故
\(\limsup a_n\le s+\varepsilon\)。第二项说明每个尾部上确界至少为
\(s-\varepsilon\)，故 \(\limsup a_n\ge s-\varepsilon\)。令
\(\varepsilon\downarrow0\) 得 \(\limsup a_n=s\)。下极限的证明把不等号和上下确界对换即可。
\end{proof}

第一项是最终上界，第二项是反复逼近。只保留第一项只能得到
\(\limsup a_n\le s\)，只保留第二项只能得到 \(\limsup a_n\ge s\)。
由同一证明还得到常用的严格版本：若 \(\limsup a_n<c\)，则最终
\(a_n<c\)；若最终 \(a_n\le c\)，则 \(\limsup a_n\le c\)。

\begin{proposition}[无穷端点]
\[
\begin{aligned}
 \limsup a_n=+\infty
 &\Longleftrightarrow
 \text{对每个 \(M\in\R\) 和每个 \(N\)，存在 \(n\ge N\) 使 \(a_n>M\)},\\
 \limsup a_n=-\infty
 &\Longleftrightarrow a_n\to-\infty.
\end{aligned}
\]
特别地，\(\limsup a_n=+\infty\) 当且仅当有子列趋于 \(+\infty\)。
下极限有完全对偶的结论。
\end{proposition}

\begin{proof}
第一行表示每个尾部都无上界，恰好等价于全部尾部上确界为
\(+\infty\)。递归地在第 \(j\) 步选 \(n_j>n_{j-1}\) 且
\(a_{n_j}>j\)，便得到趋于 \(+\infty\) 的子列；反向由该子列立即得到每个尾部无上界。

若尾部上确界 \(\beta_N\downarrow-\infty\)，给定 \(M\)，可取 \(N\)
使 \(\beta_N<M\)，于是所有 \(n\ge N\) 都有 \(a_n<M\)，即
\(a_n\to-\infty\)。反向同理。下极限结论可对 \((-a_n)\) 应用已证结果。
\end{proof}

\section{部分极限的子列刻画}


\begin{theorem}[有界数列的子列刻画]
设 \((a_n)\) 有界，记
\[
 s=\limsup a_n,\qquad \ell=\liminf a_n.
\]
则：
\begin{enumerate}
 \item 存在子列收敛到 \(s\)，也存在子列收敛到 \(\ell\)；
 \item 任意部分极限 \(L\) 都满足 \(\ell\le L\le s\)；
 \item \(s\) 是所有部分极限中的最大者，\(\ell\) 是最小者。
\end{enumerate}
\end{theorem}

\begin{proof}
设 \(\beta_n=\sup_{k\ge n}a_k\downarrow s\)。对每个 \(j\)，由上确界近似性可在尾部 \(k\ge j\) 中选指标 \(n_j\ge j\)，使
\[
 \beta_j-\frac1j<a_{n_j}\le\beta_j.
\]
为了保证严格递增，可在第 \(j\) 步把尾部起点改为 \(n_{j-1}+1\)，并使用该尾部的上确界；所得上确界仍趋于 \(s\)。夹逼给出 \(a_{n_j}\to s\)。下极限情形对 \(-a_n\) 应用同一论证。

若 \(a_{n_j}\to L\)，固定尾部起点 \(m\)。充分大时 \(n_j\ge m\)，故
\[
 \alpha_m\le a_{n_j}\le\beta_m.
\]
取 \(j\to\infty\) 得 \(\alpha_m\le L\le\beta_m\)。再令 \(m\to\infty\)，得到 \(\ell\le L\le s\)。前一步已经说明端点本身是部分极限，所以它们分别为最小和最大部分极限。
\end{proof}

对无界数列，相应端点可能为无穷。上一节的命题说明，无穷端点仍具有准确的子列含义，不需要把 \(+\infty\) 当作实数。

例如 \(a_{2n}=0\)、\(a_{2n-1}=n\) 满足
\[
 \liminf a_n=0,\qquad \limsup a_n=+\infty.
\]
上下极限允许一个端点有限、另一个端点无穷，因此比有界振荡覆盖更多发散情形。

\section{收敛判据与振荡幅度}


\begin{theorem}
实数列 \((a_n)\) 收敛到 \(L\in\R\)，当且仅当
\[
 \liminf a_n=\limsup a_n=L.
\]
\end{theorem}

\begin{proof}
若 \(a_n\to L\)，给定 \(\varepsilon>0\)，充分大时整个尾部位于
\((L-\varepsilon,L+\varepsilon)\)，故相应尾部上下确界也在闭区间
\([L-\varepsilon,L+\varepsilon]\) 中。于是上下极限均为 \(L\)。

反之，设共同值为 \(L\)。尾部下确界 \(\alpha_n\uparrow L\)，尾部上确界 \(\beta_n\downarrow L\)。给定 \(\varepsilon>0\)，取 \(N\) 使
\[
 L-\varepsilon<\alpha_N\le\beta_N<L+\varepsilon.
\]
对所有 \(n\ge N\)，有
\(\alpha_N\le a_n\le\beta_N\)，所以 \(\abs{a_n-L}<\varepsilon\)。
\end{proof}

因此振幅
\[
 \omega_n:=\sup_{k\ge n}a_k-\inf_{k\ge n}a_k
\]
对有界数列满足
\[
 a_n\text{ 收敛}\quad\Longleftrightarrow\quad\omega_n\to0.
\]
这一量度量尾部剩余的整体振荡，而不只比较相邻两项。

\section{代数运算中的等式、不等式和反例}


上极限不把加法普遍变成等式，因为两个数列的高峰可能出现在不同指标。

\begin{theorem}[和的上下极限不等式]
设 \((a_n)\)、\((b_n)\) 都有界，则
\[
 \limsup(a_n+b_n)
 \le\limsup a_n+\limsup b_n,
\]
\[
 \liminf(a_n+b_n)
 \ge\liminf a_n+\liminf b_n.
\]
若 \(a_n\to a\)，则两式中与 \(a_n\) 有关的对应关系成为等式：
\[
 \limsup(a_n+b_n)=a+\limsup b_n,
\quad
 \liminf(a_n+b_n)=a+\liminf b_n.
\]
\end{theorem}

\begin{proof}
对每个尾部，
\[
 \sup_{k\ge n}(a_k+b_k)
 \le\sup_{k\ge n}a_k+\sup_{k\ge n}b_k.
\]
令 \(n\to\infty\) 得第一式；下确界同理。若 \(a_n\to a\)，给定 \(\varepsilon>0\)，最终
\[
 a-\varepsilon\le a_n\le a+\varepsilon.
\]
于是把 \(a_n+b_n\) 夹在 \(b_n+a\pm\varepsilon\) 之间，对上下极限使用单调性，再令 \(\varepsilon\downarrow0\) 即得等式。
\end{proof}

\begin{counterexample}
令 \(a_n=(-1)^n\)、\(b_n=-(-1)^n\)。则
\[
 \limsup a_n=\limsup b_n=1,
\]
但 \(a_n+b_n=0\)，所以
\[
 \limsup(a_n+b_n)=0<2.
\]
两个高峰从不在同一指标出现。
\end{counterexample}

对非负有界数列有
\[
 \limsup(a_nb_n)\le(\limsup a_n)(\limsup b_n).
\]
若其中一列收敛到非负数，则可得到相应等式。带符号乘法需同时考虑上下极限，不能直接套用同一公式。

\begin{proposition}[标量、最大值与绝对值]
设 \((a_n)\)、\((b_n)\) 有界，\(c\in\R\)。若 \(c\ge0\)，则
\[
 \limsup(ca_n)=c\limsup a_n,\qquad
 \liminf(ca_n)=c\liminf a_n;
\]
若 \(c<0\)，则
\[
 \limsup(ca_n)=c\liminf a_n,\qquad
 \liminf(ca_n)=c\limsup a_n.
\]
此外，
\[
\begin{aligned}
 \limsup\max\{a_n,b_n\}
 &=\max\{\limsup a_n,\limsup b_n\},\\
 \liminf\min\{a_n,b_n\}
 &=\min\{\liminf a_n,\liminf b_n\},
\end{aligned}
\]
并且
\[
 \limsup\abs{a_n}
 =\max\bigl\{\abs{\liminf a_n},\abs{\limsup a_n}\bigr\}.
\]
\end{proposition}

\begin{proof}
标量公式直接来自尾部确界在正比例缩放下保持方向、在负比例缩放下交换上下确界。对最大值，每个尾部都有
\[
 \sup_{k\ge n}\max\{a_k,b_k\}
 =\max\left\{\sup_{k\ge n}a_k,\sup_{k\ge n}b_k\right\};
\]
令 \(n\to\infty\) 即得。最小值公式对偶。

记 \(\ell=\liminf a_n\)、\(s=\limsup a_n\)，并令
\(r=\limsup\abs{a_n}\)。取子列使
\(\abs{a_{n_j}}\to r\)，再由有界性抽取进一步子列
\(a_{n_{j_k}}\to L\)。于是 \(L\in[\ell,s]\)，并且
\[
 r=\abs L\le\max\{\abs\ell,\abs s\}.
\]
另一方面，分别取趋于 \(\ell,s\) 的子列并通过绝对值，说明
\(r\ge\abs\ell\) 且 \(r\ge\abs s\)。两边合并即得公式。
\end{proof}

\begin{proposition}[一个因子收敛时的乘法公式]
设 \(a_n\to a\)，而 \((b_n)\) 有界。若 \(a\ge0\)，则
\[
 \limsup(a_nb_n)=a\limsup b_n,\qquad
 \liminf(a_nb_n)=a\liminf b_n.
\]
若 \(a<0\)，则
\[
 \limsup(a_nb_n)=a\liminf b_n,\qquad
 \liminf(a_nb_n)=a\limsup b_n.
\]
\end{proposition}

\begin{proof}
先设 \(a>0\)。令 \(M\) 为 \((b_n)\) 的绝对值上界。给定
\(\varepsilon\in(0,a)\)，充分大时
\(\abs{a_n-a}<\varepsilon\)，从而
\[
 \abs{a_nb_n-ab_n}\le M\varepsilon.
\]
和的上下极限不等式应用于 \(ab_n\) 与趋零列
\((a_n-a)b_n\)，给出
\[
 \limsup(a_nb_n)=\limsup(ab_n)=a\limsup b_n,
\]
下极限同理。\(a=0\) 时由零列乘有界列得到二者都为零。
若 \(a<0\)，对 \((-a_n)b_n\) 使用正数情形，再用负号交换上下极限。
\end{proof}

两个因子都只知道上下极限时，峰值可能错位，乘法也会出现与加法相同的相关性问题。因此这些公式不能在没有收敛因子或符号条件时任意推广。

\section*{习题}
\addcontentsline{toc}{section}{习题}

\subsection*{A组　定义与基本方法}
\begin{enumerate}[label=\textbf{\thechapter.\arabic*.}]
 \exerciseitem{12.5} 对 \(a_n=(-1)^n+1/n\) 写出 \(\alpha_n,\beta_n\) 的奇偶公式，并求上下极限。
 \exerciseitem{12.6} 证明尾部上确界递减、尾部下确界递增。
 \exerciseitem{12.7} 计算下列数列的上下极限：
 \begin{enumerate}[label=(\arabic*)]
  \item \(a_n=(-1)^n n\)；
  \item \(a_n=n\sin(n\pi/2)\)；
  \item \(a_n=1+(-1)^n/n\)。
 \end{enumerate}
 \exerciseitem{12.8} 证明 \(\limsup a_n=+\infty\) 当且仅当存在子列趋于 \(+\infty\)。
\end{enumerate}

\subsection*{B组　证明与反例}
\begin{enumerate}[label=\textbf{\thechapter.\arabic*.},resume]
 \exerciseitem{12.9} 完整证明有界数列存在收敛到其上极限的子列，确保指标严格递增。
 \exerciseitem{12.10} 证明任意部分极限位于 \([\liminf a_n,\limsup a_n]\)。
 \exerciseitem{12.11} 证明 \(\limsup(-a_n)=-\liminf a_n\)。
 \exerciseitem{12.12} 证明最终 \(a_n\le b_n\) 蕴含上下极限的相应单调性。
 \exerciseitem{12.13} 证明和的上极限不等式，并给出严格不等的例子。
 \exerciseitem{12.14} 若 \(a_n\to a\)，证明
 \[
 \limsup(a_n+b_n)=a+\limsup b_n
 \]
 对有界 \((b_n)\) 成立。
 \exerciseitem{12.15} 对非负有界列证明
 \[
 \limsup(a_nb_n)\le(\limsup a_n)(\limsup b_n).
 \]
 \exerciseitem{12.16} 证明
 \[
 \limsup\max\{a_n,b_n\}
 =\max\{\limsup a_n,\limsup b_n\}.
 \]
 \exerciseitem{12.17} 对有界实数列判断 \(\limsup\abs{a_n}=\abs{\limsup a_n}\) 是否总成立，并给出正确关系。
 \exerciseitem{12.18} 证明上下极限相等且有限蕴含原列收敛。
\end{enumerate}

\subsection*{C组　综合问题}
\begin{enumerate}[label=\textbf{\thechapter.\arabic*.},resume]
 \exerciseitem{12.19} 证明有界数列收敛当且仅当尾部振幅 \(\omega_n\to0\)。
 \exerciseitem{12.21} 构造有界数列 \(a_n,b_n\)，使两个上极限都为 \(1\)，而和的上极限分别可以为 \(0,1,2\)。
 \optionalexerciseitem{12.23} 设 \(a_n>0\)。比较 \(\limsup a_{n+1}/a_n\) 与 \(a_n^{1/n}\) 的可能关系，给出至少一个不能直接交换的反例。
 \optionalexerciseitem{12.24} 证明若 \(\limsup\abs{a_n}<1\)，则存在 \(q<1\) 使 \(\abs{a_n}\le q\) 最终成立。
 \projectexerciseitem{12.26} \ 【前置：本章全部；预计用时：60分钟】设 \((a_n)\) 有界。把上、下极限写成尾部闭包交集的最大、最小点，并比较“集合表达”“尾部确界表达”“子列表达”三种证明工具。
\end{enumerate}

\section*{部分习题提示}
\begin{itemize}
 \item \exerciseref{12.5}：分别找尾部中最大的第一个偶项与最小的第一个奇项。
 \item \exerciseref{12.16}：利用 \(\max\) 与尾部上确界可交换。
 \item \exerciseref{12.17}：正确关系涉及 \(\max\{\abs{\liminf a_n},\abs{\limsup a_n}\}\)。
 \item \exerciseref{12.21}：通过让两个数列的高峰同位、错位或部分同位安排结果。
\end{itemize}
